\section{A Saúde Bucal no Brasil: Um Panorama Estrutural}

O Brasil detém um dos mais robustos sistemas de saúde pública odontológica do mundo, ancorado na Política Nacional de Saúde Bucal (PNSB) — conhecida como Brasil Sorridente —, instituída em 2004 e organicamente integrada ao Sistema Único de Saúde (SUS). Esta política representou um marco paradigmático ao transcender a histórica abordagem mutiladora e curativista, promovendo um modelo de atenção integral que abarca desde a prevenção primária até a reabilitação complexa. Contudo, a despeito dos avanços indubitáveis na redução de índices epidemiológicos como o CPO-D (Dentes Cariados, Perdidos e Obturados), o panorama atual ainda é marcado por fraturas estruturais e desafios sociodemográficos profundos.

\destaque{A desigualdade na distribuição de recursos humanos e tecnológicos constitui a principal barreira para a universalização do cuidado odontológico de excelência.}

Historicamente, a odontologia brasileira desenvolveu-se sob um modelo de mercado que favoreceu a concentração de profissionais em pólos de alta renda. Dados recentes evidenciam uma assimetria alarmante: aproximadamente 70\% dos cirurgiões-dentistas registrados no país concentram-se nas regiões Sudeste e Sul, atendendo a uma parcela demográfica que detém as melhores condições socioeconômicas e os menores índices de morbidade bucal. 

Esta concentração gera um déficit assistencial crítico em outras regiões. Na região Norte, por exemplo, a proporção chega à alarmante marca de 1 especialista em áreas complexas (como endodontia ou cirurgia bucomaxilofacial) para cada 15.000 habitantes. Tal escassez obriga pacientes a percorrerem vastas distâncias — muitas vezes dias de viagem em rotas fluviais — para acessar tratamentos especializados, o que frequentemente culmina no abandono terapêutico ou na exodontia como solução paliativa para a dor.

Ademais, a inserção da odontologia na saúde suplementar é restrita. Estima-se que apenas 30\% da população brasileira possua algum tipo de cobertura odontológica privada. Este cenário sobrecarrega o SUS, que, a despeito de sua vocação universalista, sofre com o subfinanciamento crônico. A odontologia pública, historicamente, recebe uma fração desproporcionalmente pequena do orçamento global da saúde, o que se traduz em infraestrutura defasada, escassez de insumos de alto custo e filas de espera prolongadas para procedimentos reabilitadores, como próteses e implantes.

Neste contexto de iniquidade e restrição orçamentária, a inovação tecnológica tem sido frequentemente vista como um luxo acessível apenas às clínicas de ponta dos grandes centros urbanos. Entretanto, como argumentado por Khoshfekr Rudsari et al. \cite{khoshfekr2025digital}, a digitalização na saúde possui um potencial intrínseco de democratização quando aliada a políticas públicas assertivas. A emergência dos gêmeos digitais, longe de ser um mero aprimoramento estético da prática privada, apresenta-se como uma ferramenta estratégica para subverter a lógica de exclusão, permitindo que a \textit{expertise} especializada viaje pela rede enquanto o paciente permanece em sua comunidade.
